\section{Introduction}
\label{sec:Introduction}

Deception and lying is everywhere. Courts all over the world to nothing else than to figure out the who is lying and who is telling the truth. 
When humans are actively lying, their body reacts in a certain way. Some people might look away from the person they are talking to or 
others might start to sweat or move their eyes rapidly. There are many such cues, called microexpressions that can be detected.
With the advances in computer vision models and audio detection models, we might be able to detect these small expresions and learn
about what and when persons tell the truth or when they are trying to deceive. 

The paper \cite{guo2023audio} already implemented such a architecture, that we want to build upon. They created the dataset DOLOS, as well as
using novel ways to fuse together different modalities into one trainable architecture. Our goal is to recreate the achievements of the paper,
as well as implement a new modality in the form of synthesised text. 